% =============================================================================
% MCM/ICM 论文 LaTeX 骨架模板（25 页合规版）
% 基于 2021 COMAP 官方规则 + 6 篇 O 奖论文结构提炼
% 编译：pdflatex main.tex (两遍，用于目录)
% 依赖包：geometry, fancyhdr, booktabs, siunitx, enumitem, hyperref, xcolor
% =============================================================================

\documentclass[12pt]{article}

% ---------- 页面设置：A4, 1 inch 边距 ----------
\usepackage[margin=1in,a4paper]{geometry}

% ---------- 中文字符不需要（MCM 用英文）----------
% \usepackage{ctex}  % 需要 XeLaTeX 编译

% ---------- 字体：新罗马（评委期望的学术字体）----------
\usepackage{mathptmx}
\usepackage{amsmath,amssymb,amsthm}

% ---------- 表格 ----------
\usepackage{booktabs}        % 三线表
\usepackage{multirow}        % 跨行
\usepackage{array}
\usepackage{longtable}

% ---------- 数学 ----------
\usepackage{amsmath}
\usepackage{amssymb}
\usepackage{amscd}

% ---------- 图形 ----------
\usepackage{graphicx}
\usepackage[font=small,labelfont=bf]{caption}
\usepackage{float}

% ---------- 页眉页脚（关键！匿名合规）----------
\usepackage{fancyhdr}
\pagestyle{fancy}
\fancyhf{}
\fancyhead[C]{Team \# \textbf{XXXXXXX} \quad Page \thepage \ of \pageref{LastPage}}
\fancyfoot[C]{}
\usepackage{lastpage}
\setlength{\headheight}{14pt}

% ---------- 编号 ----------
\usepackage{enumitem}

% ---------- 超链接 ----------
\usepackage[hidelinks]{hyperref}

% ---------- 颜色 ----------
\usepackage{xcolor}
\definecolor{refcolor}{RGB}{0,0,150}

% ---------- 目录深度 ----------
\setcounter{tocdepth}{3}

% ---------- 定理环境 ----------
\newtheorem{theorem}{Theorem}[section]
\newtheorem{definition}{Definition}[section]
\newtheorem{assumption}{Assumption}[section]
\newtheorem{lemma}{Lemma}[section]
\newtheorem{corollary}{Corollary}[section]

% ---------- 引用样式：[N] ----------
\usepackage[numbers,sort&compress]{natbib}
\bibliographystyle{plain}

% =============================================================================
% 主文档
% =============================================================================
\begin{document}

% ---------- 首页（无页眉）----------
\thispagestyle{plain}

\begin{center}
    \large
    \begin{tabular}{l r}
        Problem Chosen     & \textbf{X} \\
        Year               & \textbf{2024} \\
    \end{tabular}
\end{center}

\vspace{0.5cm}

\begin{center}
    {\Large \bfseries MCM/ICM}
    \vspace{0.3cm}
    {\large Summary Sheet}
    \vspace{0.3cm}
    \begin{tabular}{l r}
        Team Control Number & \textbf{XXXXXXX} \\
        Paper Title         & \textbf{Your Paper Title Here} \\
    \end{tabular}
\end{center}

\vspace{0.5cm}

\noindent
\begin{center}
    \large\bfseries Summary
\end{center}

% ---------- 摘要（关键 400-500 词）----------
% 结构：背景(1-2句) → 模型命名(1句) → Task 1-6 各1句(带数字) → 灵敏度(1句) → 结论(1句)
% 数字密度 ≥ 9 个
\begin{center}
    While [background sentence here]. This paper develops the \textbf{[Model Name, e.g., GIID Model]}
    to [core objective]. Our model has four interconnected subsystems: [subsystem 1],
    [subsystem 2], [subsystem 3], and [subsystem 4].

    \vspace{0.2cm}

    For Task 1, we [action] by [method]. The key result is that [specific number/finding, e.g.,
    \textit{``the total biomass decreases by 72.45\% under severe drought versus only 18.16\% for
    multi-species community''}].

    For Task 2, we [action] using [method]. Under [condition], [specific quantitative finding].

    For Task 3, we develop the \textit{[Self-Constructed Index Name]} to [measure what].
    Compared to [baseline], our index [specific improvement].

    For Task 4-5, sensitivity analysis shows that [parameter] is [sensitive/robust]:
    [specific number]. The model maintains [quality metric] within [range].

    For Task 6, we propose [specific recommendation] at [level]. The implementation
    results in [quantified improvement, e.g., \textit{``a 94\% reduction in [metric] over 20 years''}].

    \vspace{0.2cm}

    The main \textit{strengths} of our model are: (1) [strength 1 with mechanism],
    (2) [strength 2 with quantified advantage], and (3) [strength 3].
    The \textit{weaknesses} are: (1) [weakness 1], (2) [weakness 2], and (3) [weakness 3]---
    all of which can be addressed by [specific direction].

    \vspace{0.3cm}

    \textit{Keywords:} [keyword 1], [keyword 2], [keyword 3], [keyword 4].
\end{center}

\vfill
\pagebreak

% ---------- 目录 ----------
\tableofcontents
\clearpage

% ---------- 1 Introduction ----------
\section{Introduction}

\subsection{Background}
[1-3 段背景：从全球/宏观问题 → 具体到题目场景。每段一个概念。]

[背景第二段：题目涉及的核心矛盾。用 "However" 转折聚焦科学问题。]

\subsection{Significance}
[1-2 段：为什么这个问题重要。给出统计/经济/生态数据支撑。]

\subsection{Paper Structure}
\textbf{Note:} 这一小节是评委 3 分钟内判断"论文结构是否清晰"的关键。

In this paper, we address the problem through the following structure:
\begin{itemize}[leftmargin=*]
    \item Section 2 presents assumptions and notation.
    \item Section 3 addresses \textbf{Task 1}: [one-line summary].
    \item Section 4 addresses \textbf{Task 2}: [one-line summary].
    \item Section 5 addresses \textbf{Task 3}: [one-line summary].
    \item Section 6 addresses \textbf{Tasks 4-5}: sensitivity analysis.
    \item Section 7 addresses \textbf{Task 6}: [final discussion/recommendation].
    \item Section 8 presents model strengths and weaknesses.
\end{itemize}

\subsection{Data Description}
[数据来源、时间范围、缺失值处理、可信度说明]

\subsection{References}
[3-5 篇核心文献，标注 [1][2][3]...]

% ---------- 2 Assumptions and Notations ----------
\section{Assumptions and Notations}

\subsection{Assumptions}
\textbf{Note:} 每条假设必须带 justification——35\% 参赛队缺这个。

\begin{assumption}
[假设 1 陈述]
\end{assumption}
\noindent\textit{Justification:} [引用文献或说明推理过程，例如 "According to literature [7], ..."]

\begin{assumption}
[假设 2 陈述]
\end{assumption}
\noindent\textit{Justification:} [推理]

\begin{assumption}
[假设 3 陈述]
\end{assumption}
\noindent\textit{Justification:} [推理]

% 建议 5-7 条，不要超过 7

\subsection{Notations}
\textbf{Note:} 三列表格：符号 | 含义 | 单位。**别漏单位**。

\begin{center}
\begin{tabular}{ccc}
    \toprule
    \textbf{Symbol} & \textbf{Meaning} & \textbf{Unit} \\
    \midrule
    $N_i(t)$ & Population of species $i$ at time $t$ & individuals \\
    $B_i(t)$ & Biomass of species $i$ at time $t$ & kg/ha \\
    $W(t)$ & Soil moisture at time $t$ & mm \\
    $r_i$ & Intrinsic growth rate of species $i$ & $\mathrm{day}^{-1}$ \\
    $\alpha_{ij}$ & Competition coefficient between $i$ and $j$ & dimensionless \\
    \bottomrule
\end{tabular}
\end{center}

% ---------- 3 Task 1 ----------
\section{Task 1: [Question Description]}

\subsection{Model Development}
\textbf{Note:} 用 "we develop" 而非 "we use"——评委看到"改进"就打高分。

To answer Task 1, we develop the \textbf{[Model Name]} based on [theoretical basis].
The key modification from [classical model] is that we [specific improvement].

\subsubsection{Governing Equations}
[编号方程：(\ref{eq:xxx})]

\begin{equation}\label{eq:xxx}
    \frac{dB_i}{dt} = r_i B_i \left(1 - \frac{B_i + \sum_{j\neq i}\alpha_{ij}B_j}{K(W(t))}\right)
\end{equation}

\subsubsection{Parameters}
[参数表 + 来源]

\subsection{Results and Discussion}
\textbf{Note:} 每张图必须有正文解读段落。

Figure \ref{fig:xxx} shows [description]. From the figure, we draw the following conclusions:
\begin{itemize}[leftmargin=*]
    \item [Conclusion 1 with number]
    \item [Conclusion 2 with mechanism]
    \item [Conclusion 3 with comparison]
\end{itemize}

\begin{figure}[H]
    \centering
    \includegraphics[width=0.9\textwidth]{figures/figure1.png}
    \caption{[Concise figure caption with units]}
    \label{fig:xxx}
\end{figure}

% ---------- 4 Task 2 ----------
\section{Task 2: [Question Description]}

% ... 类似结构 ...

% ---------- 5 Task 3 (关键：自构指标) ----------
\section{Task 3: [Question Description]}

\subsection{Self-Constructed Indices}
\textbf{Note:} 自构指标是区分 M 和 O 的分水岭。至少 1 个，最好 2-3 个。

To quantitatively describe [phenomenon], we develop the following indices:

\textbf{Index 1: Total Biomass (TB)}
\begin{equation}\label{eq:tb}
    TB(t) = \sum_{i=1}^{n} B_i(t)
\end{equation}
where $B_i(t)$ is the biomass of species $i$ at time $t$.

\textbf{Index 2: [Name] Index (Abbrev)}
[公式 + 单位 + 为什么这样定义]

\subsection{Combination Analysis}
[三种组合对比 + Table]

% ---------- 6 Tasks 4-5: Sensitivity ----------
\section{Sensitivity Analysis}

\textbf{Note:} 独立章节，明确出现 "sensitivity" 一词。

\subsection{Parameter Sensitivity}
We examine the sensitivity of our model by varying [parameter] within [range]:

\begin{itemize}[leftmargin=*]
    \item \textbf{Baseline}: $\alpha = 0.5$ (default)
    \item \textbf{Low}: $\alpha = 0.1$
    \item \textbf{High}: $\alpha = 1.0$
\end{itemize}

Results in Figure \ref{fig:sensitivity} show that [trend].
\textbf{Conclusion:} The model is [robust/sensitive] with respect to $\alpha$.

\subsection{Scenario Analysis}
[多情景对比：基线 vs 极端条件]

\subsection{Initial Condition Robustness}
[初始条件扰动 + 置信区间]

% ---------- 7 Task 6: Recommendations ----------
\section{Task 6: Recommendations}

[策略一：短期措施]
\textbf{Short-term:} [specific action] which can achieve [quantified effect].

[策略二：中期措施]
\textbf{Medium-term:} [specific action] targeting [metric].

[策略三：长期措施]
\textbf{Long-term:} [specific action] with [20-year projection].

\subsection{Implementation Plan}
[时间表/路线图]

% ---------- 8 Strengths and Weaknesses ----------
\section{Strengths and Weaknesses}

\subsection{Strengths}
\textbf{Note:} 3 条，每条对应一个可命名的创新点。

\begin{enumerate}[leftmargin=*]
    \item \textbf{[Strength name]}: [mechanism + quantified advantage, e.g., "Our model reduces
    prediction error by 42\% compared to [baseline]"].
    \item \textbf{[Strength name]}: [specific methodological contribution].
    \item \textbf{[Strength name]}: [generalization/robustness claim].
\end{enumerate}

\subsection{Weaknesses}
\textbf{Note:} 3 条，每条必须给出解决方向。

\begin{enumerate}[leftmargin=*]
    \item [Weakness 1, e.g., "Parameters estimated from literature rather than field measurement."
    \textit{Direction:} Can be improved by [X]]
    \item [Weakness 2, e.g., "Limited to [scope]." \textit{Direction:} Extend to [Y]]
    \item [Weakness 3, e.g., "No feedback loop from X to Y." \textit{Direction:} Add [Z] term]
\end{enumerate}

% ---------- 参考文献 ----------
% 建议 8-12 篇，至少 3 篇来自 5 年内
\bibliographystyle{plain}
\bibliography{references}

% ---------- 附录 ----------
\appendix
\section{Derivation of [Equation]}
[详细推导]

\section{Additional Results}
[备用图、备用表]

\section{Code Availability}
[如果有代码，说明位置和语言]

% ---------- 末页 ----------
\pagebreak
\begin{center}
    \textit{End of Report}
\end{center}

\end{document}
